% Source: https://artofproblemsolving.com/community/c3h2759005_sum_of_digits_division_problem
% Stable identity: aops-topic-2759005-q01
(a) $x=1000L+100M+10N+O\le1000(L+M+N+O)$, with equality exactly when $M=N=O=0$. So the maximum is $\boxed{1000}$, attained at $1000,2000,\dots,9000$.

(b) Write $g(x)=\dfrac{1000L+100M+10N+O}{aL+M+N+O}$. A ratio of sums of positive terms lies between the smallest and largest of the individual ratios $\dfrac{1000}a$ (for $L$), $100$ (for $M$), $10$ and $1$, and it equals the largest only if every nonzero term has that ratio.
\begin{itemize}
\item If $a<10$, then $\dfrac{1000}a>100$ and the maximum $\dfrac{1000}a$ is attained only when $M=N=O=0$, so no number with $M\ne0$ maximises $g$.
\item If $a>10$, then $g(L000)=\dfrac{1000}a<100$, while $g(1900)=\dfrac{1900}{a+9}$ exceeds $\dfrac{1000}a$ (as $1900a>1000a+9000$ for $a>10$), so $L000$ does not maximise $g$.
\item If $a=10$, then $g(x)\le100$ with equality exactly when $N=O=0$; this includes $L000$ and numbers such as $1200$.
\end{itemize}
So $a=\boxed{10}$.
